\documentclass[aspectratio=169,11pt]{beamer}
\usepackage[utf8]{inputenc}
\usepackage[T1]{fontenc}
\usepackage[french]{babel}
\usepackage{helvet}
\renewcommand{\familydefault}{\sfdefault}
\usepackage{hyperref}
\definecolor{bleu}{HTML}{0000FF}
\definecolor{fond}{HTML}{F4F4EF}
\setbeamercolor{background canvas}{bg=fond}
\setbeamercolor{structure}{fg=bleu}
\setbeamercolor{title}{fg=bleu}
\setbeamercolor{frametitle}{fg=bleu,bg=fond}
\setbeamercolor{normal text}{fg=black}
\setbeamertemplate{navigation symbols}{}
\setbeamertemplate{itemize item}{\textcolor{bleu}{\tiny$\blacksquare$}}
\setbeamertemplate{footline}{\hfill\footnotesize\insertframenumber\hspace{1em}\vskip6pt}
\hypersetup{colorlinks=true,urlcolor=bleu,linkcolor=bleu}
\title{Groupe et institutions}
\subtitle{Banque africaine de développement}
\author{Veille institutionnelle}
\date{2026-11-03}

\begin{document}
\begin{frame}\titlepage\end{frame}

\begin{frame}{Groupe et institutions}
\begin{itemize}
\item Trois guichets : Banque africaine de développement, Fonds africain de développement, Fonds spécial du Nigeria
\item BAD : prêts aux pays à revenu intermédiaire et au secteur privé
\item FAD : prêts concessionnels et dons pour les pays les plus pauvres, reconstitué périodiquement
\item Fonds spécial du Nigeria : financements de projets à conditions favorables
\end{itemize}
\end{frame}

\begin{frame}{Groupe et institutions : ressources à suivre}
\small
Veille hebdomadaire.
\begin{itemize}
\item \href{https://www.afdb.org/en/about-us/african-development-fund}{Fonds africain de développement}\\{\footnotesize Page institutionnelle, publication : mensuelle}
\item \href{https://www.afdb.org/en/about-us/corporate-information/members}{Pays membres}\\{\footnotesize Liste, publication : annuelle}
\end{itemize}
\end{frame}

\end{document}
